\documentclass[../../../include/open-logic-section]{subfiles}

\begin{document}

\olfileid{sth}{choice}{woproblem}
\olsection{د ښه ترتيب مسئله}

ښکاره ده چې د ښه ترتيب منل يا نه منل ډېرې پايلې لري۔
خو د دې اصل بحث تر ډېره د هغه پر يوه همارز اصل، يعنې
د انتخاب پر اصل، راڅرخي۔ نو اوس همدغه اصل څېړو
(او همارزي به يې ثابته کړو)۔

په \citeyear{Cantor1883} کښې کانتور د ښه ترتيب
له اصله ملاتړ وکړ او هغه يې «د فکر داسې قانون چې زما
په اند بنسټيز دے، پايلې يې ډېرې دي، او په ځانګړي
ډول يې عمومي پلي کېدل د پام وړ دي» وباله
(په \citeauthor{Potter2004}
\citeyear[p.~243]{Potter2004} کښې نقل شوے)۔
خو کانتور بالاخره دې باور ته ورسېد چې دا «اصل»
ثبوت غواړي۔ د رياضي پوهانو ټولنه هم ورته نظر ته ورسېده۔

دا مسئله په \citeyear{Zermelo1904} کښې زېرملو
«هواره» کړه۔ د هغه د حل د بيانولو دپاره يو څو
تعريفونو ته اړتيا لرو۔

\begin{defn}
تابعه $f$ د \emph{انتخاب تابعه} هله او يوازې هله ده
چې $f(x) \in x$ د هر $x \in \dom{f}$ دپاره وي۔
وايو چې $f$ \emph{د $A$ دپاره د انتخاب تابعه} ده
هله او يوازې هله چې $f$ د انتخاب تابعه وي او
$\dom{f} = A \setminus \{\emptyset\}$ وي۔
\end{defn}

په وجداني توګه، د هر (ناخالي) سټ
$x \in A$ دپاره د $A$ د انتخاب تابعه
يو ځانګړے غړے، $f(x)$، له $x$ څخه
\emph{انتخابوي}۔ نو د انتخاب اصل دا دے:

\begin{axiom}[انتخاب]
هر سټ د انتخاب تابعه لري۔
\end{axiom}

زېرملو وښودله چې انتخاب ښه ترتيب لازموي، او
ښه ترتيب هم انتخاب لازموي:

\begin{thm}[په $\ZF$ کښې]\ollabel{thmwochoice}
ښه ترتيب او انتخاب همارز دي۔
\end{thm}

\begin{proof}
\emph{له چپ څخه ښي ته۔} $A$ د سټونو يو سټ
ونيسئ۔ نو $\bigcup A$ د اتحاد د اصل له مخې
شتون لري، او د ښه ترتيب له مخې داسې
$<$ شته چې $\bigcup A$ ښه ترتيبوي۔ اوس
$f(x) = \text{د $<$ له مخې تر ټولو کوچنے غړے د }x$
وټاکئ۔ دا د $A$ دپاره د انتخاب تابعه ده۔

\emph{له ښي څخه چپ ته۔} $A$ وټاکئ۔ د انتخاب
له مخې د انتخاب تابعه، $f$، د
$\Pow{A} \setminus \{\emptyset\}$
دپاره شته۔ د نامتناهيت
هاخوا بازګښت په کارولو يوه تابعه تعريف کړئ:
\begin{align*}
	g(0) &= f(A)\\
	g(\alpha) &= 
		\begin{cases}
			\text{ودرېږئ!{}} &\text{که }A = \funimage{g}{\alpha}\\
			f(A \setminus \funimage{g}{\alpha}) & \text{بل صورت کښې}\\	
		\end{cases}
\end{align*}
د «ودرېږئ!{}» اشاره يوازې د يو اوږده تعريف لنډيز
دے۔ يعنې، کله چې د لومړي ځل دپاره $A =
\funimage{g}{\alpha}$ شي، نو $g(\delta) = A$
د هر $\delta \leq \alpha$ دپاره
واخلئ۔ په لومړي سر کښې موږ يوازې دومره
باور کولے شو چې دا يو \emph{ترم} تعريفوي
(د \olref[sth][spine][recursion]{transrecursionschema}
وروسته يادونې وګورئ)؛ خو دا به وښيو چې
په رښتيا تابعه لرو۔
\textbf{د سرچينې يادښت (OLSTH-025):} د چاپي تمېدو
لارښوونه د تمېدو تر پړاو پورې ټول مخکېني ارزښتونه
هم په يوه سټ بدلوي؛ دا د لاندې يو پر يو والي له
ادعا سره نه لګېږي۔ چاپي تعريف هماغسې ساتل
شوے او دا تشه نه ده رغول شوې۔
\textbf{د سرچينې يادښت (OLSTH-026):} که ټاکلے سټ
خالي وي، د انتخاب تابعه پر هغه نه ده تعريف شوې؛
نو د چاپي بازګښت لومړنے ارزښت په دغه حالت کښې
نشته۔ بېل بنسټيز حالت نه دے وړاندې شوے۔

څرنګه چې $f$ د انتخاب تابعه ده، نو د هر
$\alpha$ دپاره (چې تعريف شوے وي)
$g(\alpha) =
f(A \setminus \funimage{g}{\alpha}) \in A \setminus
\funimage{g}{\alpha}$ لرو؛ يعنې
$g(\alpha) \notin \funimage{g}{\alpha}$۔
نو که $g(\alpha) = g(\beta)$ وي، بيا
$g(\beta) \notin
\funimage{g}{\alpha}$، يعنې
$\beta \notin \alpha$؛ او په همدې ډول
$\alpha \notin \beta$۔ نو د درې‌ګونيتوب
له مخې $\alpha = \beta$۔ په دې توګه
$g$ !!{injective} ده۔

اوس پام وکړئ چې دا به په رښتيا ودرېږي!{}؛ يعنې
داسې (تر ټولو کوچنے) ترتيبي عدد $\alpha$ شته
چې $A = g[\alpha]$ وي۔ ځکه که داسې نه وي،
نو څرنګه چې $g$ !!{injective} ده،
$\cardless{\alpha}{\Pow{A} \setminus \{\emptyset\}}$
به د هر ترتيبي عدد $\alpha$ دپاره وي، چې
له \olref[hartogs]{HartogsLemma} سره
تناقض لري۔ نو $\ran{g} = A$ هم ده۔

د دغو حقيقتونو له يو ځاے کولو سره، $g$
له يوه ترتيبي عدد څخه $A$ ته !!a{bijection}
ده۔ اوس د $g$ په وسيله $A$ ښه ترتيبولے شو۔
\end{proof}

نو ښه ترتيب او انتخاب دواړه يو ځاے ټينګېږي
يا يو ځاے لوېږي۔ خو پوښتنه لا پاتې ده: ايا
دواړه ټينګېږي که لوېږي؟

\end{document}